%% =====================================================================
%%  Nouveau programme : ce qui change (édition 2027-2028)
%%  Sources consultées le 23 octobre 2026 :
%%   - BO n° 14 du 2 avril 2026 : arrêtés du 26 février 2026 (JO du 27 mars 2026), NOR MENE2602919A (Terminale,
%%     spécialité, rentrée 2027-2028) et MENE2602917A (Première, spécialité, rentrée 2026-2027) ;
%%   - BO spécial n° 4 du 17 septembre 2026 : notes de service du 11 septembre 2026, NOR MENE2622642N (épreuve de
%%     spécialité mathématiques) et MENE2623195N (maîtrise de la langue à tous les examens à compter de la session 2027).
%% =====================================================================
\thispagestyle{otFront}
\phantomsection\pdfbookmark[0]{Nouveau programme}{bm-nouveau-programme}
\otPageTitre{Nouveau programme}{Ce qui change à la rentrée 2027, et où le trouver dans ce livre}

{\fontsize{10.5}{14.5}\selectfont
Le programme de spécialité mathématiques de Terminale publié au Bulletin officiel du 2 avril 2026 s'applique à partir de la
rentrée 2027 : il concerne les candidats au baccalauréat 2028. Il conserve l'organisation et l'essentiel des contenus du
programme de 2019. Les changements tiennent surtout à la nouvelle classe de Première, entrée en vigueur en 2026, et à la
nouvelle définition de l'épreuve écrite.\par}

\vspace{5mm}
{\fontsize{9}{11.8}\selectfont\renewcommand{\arraystretch}{1.32}\arrayrulecolor{otRule}%
\begin{tabularx}{\linewidth}{@{}>{\leavevmode\color{otNavy}\otSemi\fontsize{8.5}{10.5}\selectfont\raggedright\arraybackslash}p{30mm}
  @{\hspace{4mm}}Y@{\hspace{4mm}}>{\raggedright\arraybackslash}p{44mm}@{}}
\toprule
 & \otTete{Ce qui change} & \otTete{Dans ce livre} \\
\midrule
Fonctions sinus et cosinus & La Première ne traite plus que le cercle trigonométrique. La Terminale étudie désormais les
  fonctions $\sin$ et $\cos$ : parité, périodicité, courbes et lien avec le cercle. & Chapitre 6 : section 2 de
  « L'essentiel » réécrite, exercices 3 et 4. \\
\midrule
Plans perpendiculaires & Nouveau contenu : caractérisation par les vecteurs normaux, perpendicularité de deux plans à
  démontrer. & Chapitre 13 : théorème, démonstration et exercice 9. \\
\midrule
Ensembles et logique & Partie transversale, placée en tête du programme : notations $\varnothing$ et $\mathrm{Card}(A)$,
  négations, implications, raisonnements. & Chapitre 23, nouveau : ensembles, logique, raisonnements et rédaction. \\
\midrule
Algorithmique & Partie transversale, sans notion nouvelle : affectation notée $\leftarrow$ en langage naturel, programmation
  modulaire. & Chapitre 20 : rappels ajoutés ; chapitres 21 et 22 inchangés. \\
\midrule
Liens avec la Première & La dérivée de $x\mapsto g(ax+b)$ n'est plus vue en Première ; le produit scalaire y est étudié dans
  le plan seulement. La linéarité de l'espérance, la formule de König-Huygens et la répétition de $n\leq4$ épreuves de Bernoulli
  y sont ajoutées. & Chapitre 3 : $g(ax+b)$ traitée avec les fonctions composées. Prérequis des chapitres 3, 6, 16, 17 et 18
  mis à jour. \\
\midrule
L'épreuve écrite & Quatre heures, quatre exercices notés de $4$ à $8$ points ; deux points sur vingt pour la maîtrise de la
  langue ; moins de questionnaires sans justification. & Dix sujets type bac : exercices sur $18$ points, $2$ points de langue,
  QCM justifiés, grille officielle p.~\pageref*{grille-langue}. \\
\bottomrule
\end{tabularx}}

\vspace{4mm}
{\fontsize{9.5}{13}\selectfont\color{otMuted}
Les autres contenus (suites, limites, dérivation et convexité, continuité, logarithme, primitives et équations différentielles,
calcul intégral, géométrie dans l'espace, dénombrement, loi binomiale, sommes de variables aléatoires, loi des grands nombres)
sont repris à l'identique : les chapitres correspondants sont conservés. Les nouveautés sont signalées par la mention
\otNouveau.\par}
\clearpage

%% ---------------------------------------------------------------------
\thispagestyle{otFront}
\otPageTitre{Le programme dans ce livre}{Correspondance entre les sections du programme officiel et les chapitres}

\noindent
\begin{minipage}[t]{0.55\linewidth}\vspace{0pt}
{\fontsize{9}{11.5}\selectfont\renewcommand{\arraystretch}{1.28}\arrayrulecolor{otRule}%
\begin{tabularx}{\linewidth}{@{}Y@{\hspace{3mm}}>{\raggedleft\arraybackslash}p{17mm}@{}}
\toprule
\otTete{Section du programme} & \otTete{Chapitres} \\
\midrule
\multicolumn{2}{@{}l}{\otEtiq{otGreenDark}{Parties transversales}} \\
Vocabulaire ensembliste et logique & 23, 1, 15 \\
Algorithmique et programmation & 20 à 22 \\
\midrule
\multicolumn{2}{@{}l}{\otEtiq{otGreenDark}{Algèbre et géométrie}} \\
Combinatoire et dénombrement & 15 \\
Vecteurs, droites et plans de l'espace & 11, 12 \\
Orthogonalité et distances dans l'espace & 13 \\
Représentations paramétriques et équations cartésiennes & 13, 14 \\
\midrule
\multicolumn{2}{@{}l}{\otEtiq{otGreenDark}{Analyse}} \\
Suites & 1 \\
Limites des fonctions & 2, 10 \\
Compléments sur la dérivation & 3 \\
Continuité des fonctions d'une variable réelle & 4 \\
Fonction logarithme & 5 \\
Fonctions sinus et cosinus & 6 \\
Primitives, équations différentielles & 7, 8 \\
Calcul intégral & 9 \\
\midrule
\multicolumn{2}{@{}l}{\otEtiq{otGreenDark}{Probabilités}} \\
Succession d'épreuves indépendantes, schéma de Bernoulli & 16, 17 \\
Sommes de variables aléatoires & 18 \\
Concentration, loi des grands nombres & 19 \\
\bottomrule
\end{tabularx}}
\end{minipage}\hfill
\begin{minipage}[t]{0.4\linewidth}\vspace{0pt}
\begin{tcolorbox}[enhanced,colback=otNavyTint,colframe=otNavyTint,boxrule=0pt,arc=2mm,left=4.5mm,right=4.5mm,top=3.5mm,bottom=3.5mm,
  fontupper=\fontsize{8.6}{11.6}\selectfont\raggedright]
\otIntertitre{Textes officiels}
\textbf{Programme de Terminale}, spécialité mathématiques : arrêté du 26 février 2026, NOR MENE2602919A, Bulletin officiel
n\textsuperscript{o}\,14 du 2 avril 2026. Application : rentrée 2027.\par\vspace{2.5mm}
\textbf{Programme de Première}, spécialité mathématiques : arrêté du 26 février 2026, NOR MENE2602917A, même Bulletin.
Application : rentrée 2026.\par\vspace{2.5mm}
\textbf{Épreuve de spécialité} : note de service du 11 septembre 2026, NOR MENE2622642N, Bulletin officiel spécial
n\textsuperscript{o}\,4 du 17 septembre 2026.\par\vspace{2.5mm}
\textbf{Maîtrise de la langue} à tous les examens à compter de la session 2027 : note de service du 11 septembre 2026,
NOR MENE2623195N, même Bulletin.\par\vspace{3mm}
\href{https://www.education.gouv.fr/bo/2026/Hebdo14}{\textcolor{otNavy}{Bulletin officiel n\textsuperscript{o}\,14 du 2 avril 2026}}\par\vspace{1mm}
\href{https://www.education.gouv.fr/bo/2026/Special4}{\textcolor{otNavy}{Bulletin officiel spécial n\textsuperscript{o}\,4 du 17 septembre 2026}}\par\vspace{2.5mm}
{\color{otMuted}Textes consultés le 23 octobre 2026.}
\end{tcolorbox}
\vspace{3mm}
\begin{tcolorbox}[enhanced,colback=otGreenTint,colframe=otGreenTint,boxrule=0pt,arc=2mm,left=4.5mm,right=4.5mm,top=3.5mm,bottom=3.5mm,
  fontupper=\fontsize{8.6}{11.6}\selectfont\raggedright]
\otIntertitre{Démonstrations du programme}
Les démonstrations exigibles sont rédigées dans « L'essentiel » des chapitres concernés ou proposées en exercice, avec leur
corrigé : par exemple la limite de $(q^{n})$ (chapitre 1), la dérivée du logarithme (chapitre 5), le théorème fondamental
(chapitre 9) ou l'espérance de la loi binomiale (chapitre 17).
\end{tcolorbox}
\end{minipage}
\clearpage
